\subsection{III)}
{\noindent\large\begin{align*}\displaystyle F(x, y, z) &= \sum(0, 2, 4, 6)\end{align*}}
\subsubsection{Obtener la mínima expresión booleana}
\begin{center}
\begin{tabular}{c|ccc|c}
\rowcolor[HTML]{FFCCC9} 
\cellcolor[HTML]{9AFF99}$m$ & $x$ & $y$ & $z$ & \cellcolor[HTML]{FFFFC7}$F(x,y,z)$ \\ \hline
0                           & 0   & 0   & 0   & 1                                  \\
1                           & 0   & 0   & 1   &                                    \\
2                           & 0   & 1   & 0   & 1                                  \\
3                           & 0   & 1   & 1   &                                    \\
4                           & 1   & 0   & 0   & 1                                  \\
5                           & 1   & 0   & 1   &                                    \\
6                           & 1   & 1   & 0   & 1                                  \\
7                           & 1   & 1   & 1   &                                   
\end{tabular}
\end{center}

\begin{center}
    \begin{karnaugh-map}[2][4][1][$z$][$x\qquad y$]
        \minterms{0,2,4,6}
        \implicant{0}{4}
    \end{karnaugh-map}
\end{center}
\begin{itemize}\item En el bloque permanece constante $z$, por lo que se deja y hay que negarlo.\item La función simplificada queda:\end{itemize}
{\large\begin{align*}
    F(x,y,z) &= \overline{z}
\end{align*}}
\imagen[]{10cm}{./imagenes/ejercicio3Falstad.png}
